\begin{recipe}{Wafels}
\kicker{Dessert · zoet}
\index[register]{Wafels}
\index[register]{Eieren}
\index[register]{Boter}

\meta{35 minuten \dvd Voor 5 personen \dvd Belgisch}

\lettrine{W}{afels} maken is bij ons een zondagse bezigheid. Het beslag is in een
kwartier klaar, en als de geur door het huis trekt staat iedereen vanzelf in de keuken.

\blockrule{Ingrediënten}
\begin{ingredients}
  \ing{200 g}{roomboter}\index[register]{Boter}
  \ing{200 g}{bloem, gezeefd}
  \ing{150 g}{suiker}
  \ing{5}{eieren, gesplitst}\index[register]{Eieren}
  \ingb{zakje vanillesuiker}
  \ing{50 ml}{melk}
  \ingb{snuf zout}
\end{ingredients}

\blockrule{Bereiding}
\tip{Spatel het stijfgeklopte eiwit met een grote lepel van onder naar boven door
het beslag — zo blijven de wafels luchtig.}
\begin{steps}
  \item Klop de boter en suiker romig tot een luchtig mengsel.
  \item Splits de eieren. Voeg de dooiers één voor één toe aan het boter-suikermengsel.
  \item Roer de gezeefde bloem, het zout, de vanillesuiker en de melk door het beslag.
  \item Klop het eiwit in een aparte kom stijf.
  \item Spatel het stijfgeklopte eiwit voorzichtig door het beslag.
  \item Bak de wafels in een voorverwarmde wafelichter goudbruin.
\end{steps}
\end{recipe}
